\section{Introduction}
The micromouse project is based on the classical Micromouse competition, in which a small autonomous robot is required to solve a $16 \times 16$ maze. In the present project, however, the task was deliberately reduced to a $6 \times 6$ maze with a predefined goal position. This reduced scope was chosen to keep the project feasible within the time and budget of a university engineering project while preserving the core technical challenges of sensing, motion control, localisation, mapping, and autonomous navigation. The central task is therefore not only to move through the maze, but to do so intelligently, reliably, and efficiently within this defined academic scope.

Developing such a robot requires the integration of several engineering disciplines. The robot must be able to sense its environment, estimate relevant motion-related quantities such as position and speed, and execute controlled movements including straight-line driving, wall following, and turning manoeuvres. In addition, communication functions are required for programming, debugging, testing, and visualisation. These capabilities make the micromouse an appropriate platform for combining embedded programming, electronics, control engineering, and navigation algorithms in a single system.

A key aspect of the project is the use of closed-loop control. Since the robot operates in a constrained maze environment, open-loop motion alone is not sufficient to achieve reliable and repeatable navigation. Sensor feedback must be processed continuously in order to correct deviations in motion and to maintain accurate alignment within the maze. This makes control design a fundamental part of the overall system.

From a hardware perspective, the project is centered around a microcontroller-based embedded system. The required functionality includes digital input and output, interrupts, pulse-width modulation, serial communication, quadrature encoder interfaces, and analogue-to-digital conversion. These peripheral features provide the basis for connecting sensors, actuators, and debugging interfaces, and for implementing the software library needed to operate the robot.

Overall, the micromouse project serves as a practical case study in the development of an autonomous robotic system. It combines the design of hardware and software with the implementation of sensing, control, and mapping strategies. The result is a compact but technically demanding platform that demonstrates how embedded systems and control methods can be applied to autonomous navigation in a structured environment.

\subsection{Aim}
The main aim of this project is to design, implement, and evaluate a functioning micromouse robot that is capable of navigating autonomously through a reduced $6 \times 6$ maze environment.

\subsection{Objectives}
The objectives of this project are as follows:
\begin{itemize}
    \item To develop a hardware platform including microcontroller, sensors, actuators, and power supply.
    \item To implement the required low-level embedded functionality for interfacing with the robot hardware.
    \item To design and integrate software modules for sensing, motion control, communication, and navigation.
    \item To implement closed-loop motor and motion control strategies for accurate movement in the maze.
    \item To test and evaluate the performance of the sensor system, controller, and navigation behaviour.
    \item To analyse the limitations of the implemented system and identify possible improvements.
\end{itemize}
